\documentclass[10pt,a4paper]{article}
\usepackage[T1]{fontenc}
\usepackage{graphicx}
\usepackage{mathtools}
\usepackage{amssymb}
\usepackage{amsthm}
\usepackage{thmtools}
\usepackage{xcolor}
\usepackage{nameref}
\usepackage{float}
\usepackage{booktabs}

%下面的是配置是可以写中文注释的python环境:
\usepackage{listings}
\usepackage{color}
\definecolor{dkgreen}{rgb}{0,0.6,0}
\definecolor{gray}{rgb}{0.5,0.5,0.5}
\definecolor{mauve}{rgb}{0.58,0,0.82}
\lstset{frame=tb,
	language=Python,
	aboveskip=3mm,
	belowskip=3mm,
	showstringspaces=false,
	columns=flexible,
	basicstyle={\small\ttfamily},
	numbers=left,%设置行号位置none不显示行号
	%numberstyle=\tiny\courier, %设置行号大小  
	numberstyle=\tiny\color{gray},
	keywordstyle=\color{blue},
	commentstyle=\color{dkgreen},
	stringstyle=\color{mauve},
	breaklines=true,
	breakatwhitespace=true,
	escapeinside=``,%逃逸字符(1左面的键)，用于显示中文例如在代码中`中文...`
	tabsize=4,
	extendedchars=false} %解决代码跨页时，章节标题，页眉等汉字不显示的问题  
\usepackage{ctex}
\begin{document}
		\section*{\centering Appendix2：Minute-frequency Trading}
	\subsection*{1. database and table}
	\begin{lstlisting}
		/* 1. 存放处理后的数据，例如k线、个股收益 */
		create database "dfs://G3XuZhitong_cookedDB"
		partitioned by VALUE(2023.08M..2023.09M), HASH([SYMBOL,256])
		engine = "TSDB"
		
		/* 1.1 分钟K线表 */
		create table "dfs://G3XuZhitong_cookedDB".'min_k_line'(
		TradeDate DATE,
		SecurityID SYMBOL,
		MinuteStart MINUTE,
		Open DOUBLE,
		Close DOUBLE,
		High DOUBLE,
		Low DOUBLE,
		Volume DOUBLE,
		TradeValue DOUBLE
		)
		partitioned by TradeDate, SecurityID
		sortColumns = ['SecurityID','TradeDate','MinuteStart']
		keepDuplicates=ALL
		
		
		/* 2. 存放计算得到的因子数据 */
		factornames = ['SVOL_10','DSVOL_10']
		create database "dfs://G3XuZhitong_FactorDB"
		partitioned by VALUE(factornames), VALUE(2023.08M..2023.09M)
		engine = "TSDB"
		
		/* 2.1 因子原始值表 */
		create table "dfs://G3XuZhitong_FactorDB".'factor_minute'(
		TradeDate DATE,
		SecurityID SYMBOL,
		MinuteStart MINUTE,
		factorname STRING,
		value DOUBLE
		)
		partitioned by factorname, TradeDate
		sortColumns = ['SecurityID','TradeDate','MinuteStart']
		keepDuplicates=ALL
		
		/* 2.2 因子暴露表 */
		create table "dfs://G3XuZhitong_FactorDB"."exposure_minute"(
		TradeDate DATE,
		SecurityID SYMBOL,
		MinuteStart MINUTE,
		factorname STRING,
		exposure DOUBLE,
		neutral_type STRING
		)
		partitioned by factorname, TradeDate
		sortColumns = ['SecurityID','TradeDate','MinuteStart']
		keepDuplicates=ALL
		
		/* 2.3 因子表现表 */
		create table "dfs://G3XuZhitong_FactorDB"."performance_minute"(
		factorname STRING,
		neutral_type STRING,
		startDate DATE,
		endDate DATE,
		holding_minutes INT,
		nGroups INT,
		weightMethod STRING,
		costrate DOUBLE,
		rf DOUBLE,
		IC DOUBLE,
		IRIC DOUBLE,
		rankIC DOUBLE,
		rankIRIC DOUBLE,
		TotalReturn DOUBLE,
		AnnualizedReturn DOUBLE,
		AnnualizedVolatility DOUBLE,
		SharpeRatio DOUBLE,
		annSharpeRatio DOUBLE,
		MaxDrawdown DOUBLE,
		WinRate DOUBLE
		)
		sortColumns = ['factorname','neutral_type','startDate','endDate']
	\end{lstlisting}
	
	\subsection*{2. generate minute K-lines}
	\begin{lstlisting}
		cuhkdb = database("dfs://cuhkDB")
		filtered_data =
		SELECT 
		TradeDate,
		TradeTime,
		SecurityID,
		minute(TradeTime) AS MinuteStart, 
		LastPrice,
		TotalVolumeTrade,
		TotalValueTrade
		FROM cuhkdb.cuhk_snapshot
		WHERE TradingPhaseCode=="TRADE"
		order by TradeDate, TradeTime, SecurityID
		
		cookeddb = database("dfs://G3XuZhitong_cookedDB")
		cookeddb.min_k_line.append!(
		SELECT TradeDate,
		SecurityID,
		MinuteStart,
		first(LastPrice) AS Open,
		last(LastPrice) AS Close, 
		max(LastPrice) AS High,
		min(LastPrice) AS Low,
		last(TotalVolumeTrade)-first(TotalVolumeTrade) AS Volume,
		last(TotalValueTrade)-first(TotalValueTrade) AS TradeValue
		FROM filtered_data
		GROUP BY TradeDate, SecurityID, MinuteStart
		ORDER BY TradeDate, SecurityID, MinuteStart
		)
		
	\end{lstlisting}
	
	\subsection*{3.load factor}
	\begin{lstlisting}
		pricedata = loadTable("dfs://G3XuZhitong_cookedDB", "min_k_line")
		
		factorload = call(VOL_CF, pricedata, 10)
		if(factorload.factorname[0] in (select distinct factorname from loadTable("dfs://G3XuZhitong_FactorDB", "factor_minute"))){
			print("因子 " + factorload.factorname[0] + " 已存在于因子表中，跳过添加。")
		}
		else{
			upsert!(loadTable("dfs://G3XuZhitong_FactorDB", "factor_minute"), factorload, KeyColNames=`TradeDate`SecurityID`MinuteStart`factorname)
			print("因子 " + factorload.factorname[0] + " 已成功添加到因子表中。")
		}
	\end{lstlisting}
	
	\subsection*{4. process factor}
	\begin{lstlisting}
		/* 对单个因子、单个截面做处理,分位数0.01-0.99去极值，zscore标准化 */
		def processFactor(mutable valueVec, winsorizeLower=1, winsorizeUpper=99,neutral_type="winsorized_zscore") {
			/* 1. 去空值 */
			valueVec_mask = isNull(valueVec)
			valueVec = valueVec[not valueVec_mask]
			
			if(neutral_type == "winsorized_zscore") {
				/* 2. 去极值 */
				lower, upper = percentile(valueVec, winsorizeLower), percentile(valueVec, winsorizeUpper)
				valueVec = clip(valueVec, lower, upper)
				
				/* 3. Z-score */
				mu = avg(valueVec)
				sigma = std(valueVec)
				if(sigma == 0.0) return take(0.0, size(valueVec))
				return (valueVec - mu) / sigma 
			}
		}
		
		/* 将因子暴露存入exposure_minute表 */
		def saveFactorExposureminute(fname, type="winsorized_zscore") 
		{
			/* 读取因子数据 */
			factorData = select *, type as neutral_type from database("dfs://G3XuZhitong_FactorDB").factor_minute 
			where factorname = fname and value is not null
			
			/* 按TradeDate,MinuteStart分组处理 */
			processedData = select 
			TradeDate,
			SecurityID,
			MinuteStart,
			factorname,
			processFactor(value) as exposure,
			neutral_type
			from factorData
			context by TradeDate, MinuteStart
			
			/* 保存结果到exposure_minute表 */
			upsert!(database("dfs://G3XuZhitong_FactorDB").exposure_minute,(
			select 
			TradeDate,
			SecurityID,
			MinuteStart,
			factorname,
			exposure,
			neutral_type
			from processedData
			))
			
			print("已处理新因子：" + fname + ",type:" + type)
		}
		
		/* 将未按neutral_type处理的因子值处理并存入exposure_minute表 */
		factor_minute   = loadTable("dfs://G3XuZhitong_FactorDB", "factor_minute")
		exposure_minute = loadTable("dfs://G3XuZhitong_FactorDB", "exposure_minute")
		
		fm_set = select distinct factorname from factor_minute 
		ep_set = select distinct factorname from exposure_minute where neutral_type = "winsorized_zscore"
		missing_factors = fm_set.factorname[not fm_set.factorname in ep_set.factorname]
		if(size(missing_factors) > 0){
			print("待处理因子列表：")
			print(missing_factors)
			for(fname in missing_factors){
				saveFactorExposureminute(fname, "winsorized_zscore")
			}
		}
		else{
			print("无待处理因子")
		}
	\end{lstlisting}
	
	\subsection*{5. back testing}
	\begin{lstlisting}
		// -------------------------------
		fname              = "VOL_CF_10min_8delta"        // 因子名
		type               = "winsorized_zscore"        // 处理类型
		holding_minutes    = 10     // 持仓分钟
		nGroups            = 5                    // 分组数
		weightMethod       = "equal"                // 权重方案： "rank" / "equal" / "exposure"
		costrate           = 0.000                    // 单边手续费
		universeFilterFunc = NULL                  // 可选：自定义函数筛选股票（NULL 表示不筛选）
		rf = 0.03                                  // 无风险利率
		startDate          = 2023.08.14            // 回测开始日期
		endDate            = 2023.09.14            // 回测结束日期
		// -------------------------------
		
		/* 载入因子表和分钟k线 */
		exposureTbl = loadTable("dfs://G3XuZhitong_FactorDB", "exposure_minute")   
		priceTbl    = loadTable("dfs://G3XuZhitong_cookedDB", "min_k_line")
		
		/* 取出目标因子并生成买入时间entry_dt,卖出时间exit_dt */
		exposure = select TradeDate, SecurityID, MinuteStart, factorname,exposure, neutral_type 
		from exposureTbl 
		where TradeDate between startDate and endDate and factorname = fname and neutral_type = type   // 根据 factorName 过滤
		
		exposure = select *, concatDateTime(TradeDate, second(time(MinuteStart))) as entry_dt from exposure
		exposure = select *, concatDateTime(TradeDate, second(time(MinuteStart+holding_minutes-1))+59) as exit_dt from exposure
		// 过滤掉开盘前、收盘后无法交易的时点
		exposure = select * from exposure where minute(exit_dt) <= 14:59m  and minute(entry_dt) >= 09:30m
		// 补充k线每一窗口的开始、结束时点
		price = select *, concatDateTime(TradeDate, second(time(MinuteStart))) as Open_dt, concatDateTime(TradeDate, second(time(MinuteStart+1))-1) as Close_dt from priceTbl
		// 因子表连接价格表，获取买入价和卖出价
		// entry price join
		entryJoin = select a.TradeDate, a.SecurityID, a.entry_dt, a.exit_dt, a.exposure, a.neutral_type
		, p.Open  as entry_open
		from exposure  a
		left join price  p on a.TradeDate = p.TradeDate and a.SecurityID = p.SecurityID and a.entry_dt = p.Open_dt
		// exit price join
		fullJoin = select j.*, p.Close as exit_close
		from entryJoin  j
		left join price  p on j.TradeDate = p.TradeDate and j.SecurityID = p.SecurityID and j.exit_dt = p.Close_dt
		// 计算向前收益
		fullJoin = select *, 
		iif(isNull(entry_open) or isNull(exit_close) or entry_open==0, NULL, (exit_close/entry_open - 1.0)) as fwdReturn
		from fullJoin
		
		/* IC */
		ICTbl = select entry_dt, 
		corr(exposure, fwdReturn)        as IC, 
		corr(rank(exposure), rank(fwdReturn)) as rankIC,
		count(exposure) as N
		from fullJoin
		group by entry_dt
		// IC统计数据
		ICStats = select 
		avg(IC)       as meanIC,
		std(IC)     as stdIC,
		skew(IC)    as skewIC,
		kurtosis(IC) as kurtosisIC,
		avg(IC) / std(IC) * sqrt(count(IC)) as tStatIC,
		avg(IC) / std(IC) as irIC,
		avg(rankIC)      as meanrankIC,
		std(rankIC)    as stdrankIC,
		skew(rankIC)   as skewrankIC,
		kurtosis(rankIC) as kurtosisrankIC,
		avg(rankIC) / std(rankIC) * sqrt(count(rankIC)) as tStatrankIC,
		avg(rankIC) / std(rankIC) as irrankIC,
		count(entry_dt)      as numPeriods
		from ICTbl
		
		ICmain = select  meanIC as IC, irIC as IRIC, meanrankIC as rankIC, irrankIC as rankIRIC from ICStats
		
		/* 分组收益率计算 */
		// 将每期股票按因子值分组
		grouped = select TradeDate, SecurityID, entry_dt, exposure, fwdReturn, rank(exposure,groupNum=nGroups) as group 
		from fullJoin
		context by entry_dt
		// 每期的组内收益率统计
		groupedSummary = select entry_dt, group, 
		avg(fwdReturn) as meanReturn, 
		std(fwdReturn) as stdReturn, 
		count(fwdReturn) as N
		from grouped
		group by entry_dt, group
		order by entry_dt, group
		// 各组收益率统计         
		groupedStats = select 
		group,
		avg(meanReturn) as avgReturnPerGroup,
		std(meanReturn) as stdReturnPerGroup
		from groupedSummary
		group by group
		
		/* 计算多空组合收益:long-short等权重 */
		if(ICStats.meanIC > 0){
			lsSpread = select t.entry_dt, (t.meanReturn - b.meanReturn) as longShort
			from (select entry_dt, meanReturn, group from groupedSummary where group = nGroups-1) t
			join (select entry_dt, meanReturn, group from groupedSummary where group = 1) b
			on t.entry_dt = b.entry_dt
		}
		else{
			lsSpread = select t.entry_dt, (t.meanReturn - b.meanReturn) as longShort
			from (select entry_dt, meanReturn, group from groupedSummary where group = 1) t
			join (select entry_dt, meanReturn, group from groupedSummary where group = nGroups-1) b
			on t.entry_dt = b.entry_dt
		}
		
		/* PNL */
		// 根据 holding_minutes 进行收益率采样
		if(holding_minutes > 1){
			PNL = select * from
			(select *,row_number() over (order by entry_dt) as seq
			from lsSpread )
			where seq % holding_minutes ==1
		}
		else{
			PNL = select * from
			(select *,row_number() over (order by entry_dt) as seq
			from lsSpread )
			
		}
		// 减去交易成本
		update PNL set longShort = longShort - 2 * costrate
		// 计算累计收益率和净值
		PNL = select *, cumprod(1+longShort) - 1 as netret from PNL
		
		/* 绩效指标计算 */
		totalperiods = size(PNL.entry_dt)         // 回测的总 period 数
		periodsPerDayTbl = select date(entry_dt), count(*) as cnt from PNL group by date(entry_dt) // 每天的 period 数
		avgPeriodsPerDay = avg(periodsPerDayTbl.cnt)  // 平均每天调仓次数
		tradingDaysPerYear = 252.0                    // 一年交易日天数
		periodsPerYear = avgPeriodsPerDay * tradingDaysPerYear // 一年 period 数
		// 1. 最大回撤
		nr = PNL.netret 
		runningMax = nr.cummax()    // 向量的累积最大值（每点之前的最高净值）
		drawdown = nr / runningMax - 1.0   // 每点的回撤
		maxDrawdown = min(drawdown)        // 最大回撤
		// 最大回撤发生时间点
		idxMaxDD = at(drawdown == maxDrawdown)
		maxDD_time = iif(size(idxMaxDD)>0, PNL.entry_dt[idxMaxDD], NULL)
		// 2. 年化
		meanPeriodReturn = avg(PNL.longShort) // 每期平均收益
		stdPeriodReturn = std(PNL.longShort)  // 单期收益标准差
		totalreturn = PNL.netret[totalperiods-1] // 回测时间段总收益率
		annVol = pow(pow(1+meanPeriodReturn,2)+pow(stdPeriodReturn,2), periodsPerYear)-pow(pow(1+meanPeriodReturn,2),periodsPerYear) // 年化波动率 
		annReturn = pow(1 + totalreturn, periodsPerYear/totalperiods) - 1.0 // 年化收益率
		// 3. 夏普比率
		rfperiod = pow(1 + rf, 1.0/periodsPerYear) - 1.0 // 每期无风险收益率
		sharpe = (meanPeriodReturn - rfperiod) / stdPeriodReturn
		annsharpe = (annReturn-rf)/annVol //年化夏普
		// 4. 胜率
		winRate = sum(PNL.longShort > 0) * 1.0 / totalperiods
		// 输出绩效指标
		performance = table(  totalreturn as TotalReturn,
		annReturn as AnnualizedReturn,
		annVol as AnnualizedVolatility,
		sharpe as SharpeRatio,
		annsharpe as annSharpeRatio,
		maxDrawdown as MaxDrawdown,
		winRate as WinRate )
		
		/* 绘图 */
		plot(PNL.netret, PNL.entry_dt)
		// plot(ICTbl.spearmanIC, ICTbl.entry_dt)
		
		
		/* 保存因子绩效到因子表现表 */
		upsert!(loadTable("dfs://G3XuZhitong_FactorDB", "performance_minute"), 
		(select fname as factorname,
		type as neutral_type,
		startDate as startDate,
		endDate as endDate,
		holding_minutes as holding_minutes,
		nGroups as nGroups,
		weightMethod as weightMethod,
		costrate as costrate,
		rf as rf,
		ICmain.IC as IC,
		ICmain.IRIC as IRIC,
		ICmain.rankIC as rankIC,
		ICmain.rankIRIC as rankIRIC,
		performance.TotalReturn as TotalReturn,
		performance.AnnualizedReturn as AnnualizedReturn,
		performance.AnnualizedVolatility as AnnualizedVolatility,
		performance.SharpeRatio as SharpeRatio,
		performance.annSharpeRatio as annSharpeRatio,
		performance.MaxDrawdown as MaxDrawdown,
		performance.WinRate as WinRate
		from performance, ICmain),
		KeyColNames=['factorname','neutral_type','startDate','endDate','holding_minutes','nGroups','weightMethod','costrate','rf']
		)
	\end{lstlisting}
\end{document}